\newcommand{\ch}{\mathfrak{X}}

From now on, let $G$ be a (connected) reductive group. We also fix a Borel subgroup $B \subseteq G$ and a maximal torus $T \subseteq B$.

\begin{definition}\label[definition]{def:sph}
	A \emph{spherical $G$-variety} $X$ is a (normal, irreducible) $G$-variety with an open $B$-orbit.
\end{definition}
\begin{remark*}
	Since $X$ has an open $B$-orbit, it has an open $G$-orbit.
\end{remark*}

\begin{example}
	When $G = B = T$ (and we may further require the $T$-action is faithful), the spherical variety $X$ is a \emph{toric variety}.
\end{example}

We write $\ch(A)$ for the \emph{character lattice} of an algebraic group $A$. For a reductive group $G$, the character lattice $\ch(B) \cong \ch(T)$, and they are isomorphic to the lattice $\mathbb{Z}^r$, where $r := \rank(G)$ is the dimension of $T$.




% The representation theory of reductive groups is well-studied.
% \begin{lemma}[c.f., e.g., {\cite[\S15.2]{MT_Lin_alg_grp_fin_grp_Lie}}]
% 	Every finite-dimensional $G$-module is semisimple (i.e., direct sum of simple $G$-modules). For every finite-dimensional simple $G$-module $V$, there exists a unique (up to scaling) non-zero vector $v \in V$ that satisfies 
% 	\begin{equation*}
% 		b \cdot v = \lambda(b) v,
% 	\end{equation*}
% 	for some $\lambda \in \ch(B)$ and every $b \in B$. The unique $\lambda$ that satisfies the equation is called the \emph{highest weight} of $V$. Moreover, two finite-dimensional simple $G$-modules are isomorphic iff they have the same highest weight.
% \end{lemma}

We will use the following notation. Given a finite-dimensional $G$-module $V$, for every character $\lambda \in \ch(B)$,
\begin{itemize}
	\item $V_\lambda$, the \emph{$\lambda$-isotypic subspace}, is the sum of all (isomorphic) simple $G$-submodules of $V$ with the highest weight $\lambda$; since the vector space $V$ is finite-dimensional, we have $\bigoplus_{\lambda \in \ch(B)} V_\lambda = V$;

	\item $V_\lambda^{(B)}$ is the subset of $0$ and \emph{$B$-semi-invariant vectors} in $V_\lambda$, i.e., $V_\lambda^{(B)} = \{v \in V \mid b \cdot v = \lambda(b) v, \forall b \in B\}$; the subset $V_\lambda^{(B)}$ is a vector space, and it is of dimension $1$ iff $V_\lambda$ is simple.
\end{itemize}

For a general reference about the representation theory of reductive groups, see, for example, \cite{MT_Lin_alg_grp_fin_grp_Lie}. 

Given a $G$-variety $X$, the function field $\mathbb{C}(X)$ is an infinite-dimensional $G$-module, but for every $\lambda \in \ch(B)$, we can still define 
\begin{equation*}
	\mathbb{C}(X)_\lambda^{(B)} := \{f \in \mathbb{C}(X) \mid b \cdot f = \lambda(b) f, \forall b \in B\}.
\end{equation*}
And similar to toric varieties, we can define 
\begin{equation*}
	M := \{\lambda \in \mathfrak{X}(B) \mid \mathbb{C}(X)_\lambda^{(B)} \neq 0 \},
\end{equation*}
which is a sub-lattice of $\mathfrak{X}(B)$; its rank is called the \emph{rank} of $X$.

Given a spherical variety $X$, any $B$-invariant rational function on $X$ is constant; equivalently, for every character $\lambda \in M$, we have $\dim(\mathbb{C}(X)_\lambda^{(B)}) = 1$. For every function $f_\lambda \in \mathbb{C}(X)_\lambda^{(B)} \setminus \{0\}$, the divisor $\div(f_\lambda)$ is principal and $B$-invariant; conversely, any principal $B$-invariant divisor is determined by some character $\lambda \in M$.

\subsubsection{\textit{G}-invariant Valuations}\label{sec:prelim:sph:val}
\newcommand{\V}{\mathcal{V}}

Similar to toric varieties, let $N := \Hom(M, \mathbb{Z})$, $N_\mathbb{Q} := \Hom(M, \mathbb{Q})$, and we denote the perfect pairing by $\<\cdot, \cdot\> : M_\mathbb{Q} \times N_\mathbb{Q} \to \mathbb{Q}$.

Given a valuation $v : \mathbb{C}(X)^\times \to \mathbb{Q}$, we can define a map $\rho(v)(\lambda) := v(f_\lambda)$ for every character $\lambda \in M$; we have $\rho(v) \in N_\mathbb{Q}$, i.e.,
\begin{equation}
\begin{aligned}
	\rho: \{\text{valuations of }\mathbb{C}(X)\} \to N_\mathbb{Q}
\end{aligned}
\end{equation}
is well-defined.

% valuation cone
\begin{lemdefn}[c.f., e.g., {\cite[\S4]{Brion_var_sph}}]\label[lemdefn]{def:valcone}
	After restriction,
	\begin{equation*}
		\rho : \{G\text{-invariant valuations of }\mathbb{C}(X)\} \to N_\mathbb{Q}
	\end{equation*}
	is injective. Moreover, its image is a co-simplicial convex cone, called \emph{valuation cone} of $X$, denoted by $\V$.
\end{lemdefn}
\begin{remark*}
	In other words, each $G$-invariant valuation of $\mathbb{C}(X)$ is uniquely determined by its values on $B$-semi-invariant rational functions, or equivalently, principal $B$-invariant divisors.
\end{remark*}

As a result, we identify a $G$-invariant valuation $v$ with a vector $\rho(v) \in \V$.

As the cone $\V$ is polyhedral, the following lemma is intuitive.
\begin{lemma}\label[lemma]{lem:sph:DivValEQVal}
	Assume $X$ is a projective spherical $G$-variety, we have $\Val^G_X = \DVal^G_X$.
\end{lemma}
\begin{remark*}
	Recall that, in this paper we assumed valuations are $\mathbb{Q}$-valued. Note that there always exists $G$-invariant $\mathbb{R}$-valuations which are not of the form $c \: \ord_E$ for some prime divisor $E$ over $X$ and $c \in \mathbb{R}_{\geq 0}$; for instance, quasi-monomial ($\mathbb{R}$-)valuations (see, for example, \cite[\S3]{JM_val_asym_inva_seq_i}).
\end{remark*}
\begin{sketchproof}
	According to \cite[Proposition 2, \S2.4]{Brion_var_sph}, there is a toroidal $G$-equivariant birational morphism $\pi: X' \to X$. For every $G$-invariant valuation $v \in \Val^G_X$, there is a subdivision of $\V$ where $\mathbb{Q}_{\geq 0} v \subseteq \V$ is a cone. The subdivision of $\V$ corresponds to birational model $X''$ of $X'$ (thus a birational model of $X$), and the ray $\mathbb{Q}_{\geq 0} v$ defines a prime divisor of $X''$.
\end{sketchproof}

This lemma is stronger than \cite[Proposition 2.19]{Golota_Delta_Fano_large_aut}, which proves the subset $\DVal^G_X$ is dense in the space of real-valued $G$-invariant valuations.



